% ============================================================
%  HSE University — Beamer Presentation Template
%  Resolution  : 1920 × 1080  (aspectratio=169)
%  Compiler    : pdfLaTeX  (Overleaf)
%  Theme file  : hse.sty  (must be in the same folder)
% ============================================================
\documentclass[aspectratio=169,12pt,t]{beamer}

% ── Russian language support ─────────────────────────────────
\usepackage{cmap}
\usepackage[T2A]{fontenc}
\usepackage[utf8]{inputenc}
\usepackage[english,russian]{babel}

% ── HSE theme ────────────────────────────────────────────────
\usepackage{hse}

\setbeamertemplate{headline}[hse theme]
\setbeamertemplate{frametitle}[hse theme]
\setbeamertemplate{footline}[hse theme]

% ── Override localisation macros (now UTF-8 is active) ───────
\renewcommand{\LogoPlaceholderText}{ВШЭ\,/\,HSE\\LOGO}
\renewcommand{\ThankYouText}{Спасибо!}
\renewcommand{\ContactsText}{Контакты}
\renewcommand{\FeedbackText}{Отзыв о докладе}

% ── Additional packages ───────────────────────────────────────
\usepackage{qrcode}   % QR codes on the final slide

% ============================================================
%  PRESENTATION METADATA — fill in before use
% ============================================================
\newcommand{\EventName}{Название конференции}
\newcommand{\EventLocation}{Москва, Россия}
\newcommand{\PresentationTitle}{Заголовок презентации}
\newcommand{\PresentationSubtitle}{Подзаголовок или краткое описание темы}
\newcommand{\AuthorName}{Иван Иванович Иванов}
\newcommand{\AuthorRole}{Доцент, Факультет компьютерных наук}
\newcommand{\AuthorAffiliation}{НИУ ВШЭ — Высшая школа экономики}
\newcommand{\PresentationDate}{16 апреля 2026}

% Text shown in the footer (left side)
\renewcommand{\FooterLeft}{\AuthorName\ \textbullet\ \AuthorAffiliation}

% URLs for the final slide QR codes
\newcommand{\ContactURL}{https://example.com/contact}
\newcommand{\FeedbackURL}{https://forms.gle/example}

% ── Logo ─────────────────────────────────────────────────────
% To use a real logo, uncomment and edit the line below,
% then comment out or redefine \logobox in hse.sty:
%
%   \renewcommand{\logobox}{\includegraphics[height=1.2cm]{hse-logo}}
%
% The logo file (e.g. hse-logo.pdf) must be in the same folder.

% ============================================================
\begin{document}

% ── TITLE SLIDE ──────────────────────────────────────────────
\begin{frame}[plain,noframenumbering]
  \titlepage
\end{frame}

% ── TABLE OF CONTENTS ────────────────────────────────────────
% [nonavbar] keeps this slide out of the nav-bar dots
\begin{frame}[nonavbar]{Содержание}
  \tableofcontents
\end{frame}

% ============================================================
\section{Введение}
% ============================================================

\begin{frame}{Обычный слайд}
  Это стандартный слайд. Навигационная шапка вверху показывает
  все разделы; активный раздел выделен жёлтой плашкой.
  Заголовок фрейма отделён от шапки тонкой жёлтой полосой.

  \begin{itemize}
    \item Пункт списка первый
    \item Пункт списка второй
    \item Пункт списка третий
  \end{itemize}
\end{frame}

\begin{frame}{Выделения и акцентные блоки}
  Встроенное красное выделение работает в обычном тексте:
  \emred{это важный фрагмент} и продолжение предложения.

  В формулах тоже работает:
  \[
    E = \emred{mc^2}, \qquad F = \emred{ma}
  \]

  \begin{emphblock}{Ключевой результат}
    Самостоятельный блок с красной рамкой — для определений,
    теорем, главных выводов или предупреждений.
  \end{emphblock}
\end{frame}

\begin{frame}{Стандартные блоки Beamer}
  \begin{block}{Обычный блок}
    Используется для нейтральных пояснений.
  \end{block}

  \begin{alertblock}{Предупреждение}
    Используется для ошибок или критически важных замечаний.
  \end{alertblock}

  \begin{exampleblock}{Пример}
    Используется для иллюстраций или демонстраций.
  \end{exampleblock}
\end{frame}

% ── Section divider ──────────────────────────────────────────
\sectionslide{Методология}

% ============================================================
\section{Методология}
% ============================================================

\begin{frame}{Одно изображение по центру}
  % \oneimage{файл}{подпись}{ширина}
  \oneimage{example-image-a}{Рисунок 1. Описание изображения}{0.58\textwidth}
\end{frame}

\begin{frame}{Два изображения рядом}
  % \twoimages{лево}{подпись лево}{право}{подпись право}
  \twoimages
    {example-image-a}{Рисунок 1. Левая подпись}
    {example-image-b}{Рисунок 2. Правая подпись}
\end{frame}

\begin{frame}{Сетка 2×2}
  % \fourgrid{img1}{cap1}{img2}{cap2}{img3}{cap3}{img4}{cap4}
  \fourgrid
    {example-image-a}{Верхний левый}
    {example-image-b}{Верхний правый}
    {example-image-c}{Нижний левый}
    {example-image}  {Нижний правый}
\end{frame}

\begin{frame}{Изображение и текст}
  % \imagewithtext{файл}{содержимое правой колонки}
  \imagewithtext{example-image-a}{%
    \textbf{Ключевые наблюдения}
    \begin{itemize}
      \item Первое важное наблюдение из рисунка
      \item Второй вывод с пояснением
      \item Перспективы дальнейшего анализа
    \end{itemize}
  }
\end{frame}

% ── Section divider ──────────────────────────────────────────
\sectionslide{Результаты}

% ============================================================
\section{Результаты}
% ============================================================

\begin{frame}{Два столбца}
  % \twocols{левый контент}{правый контент}
  \twocols{%
    \textbf{Теоретическая часть}\\[6pt]
    Произвольный контент: текст,
    формулы, списки.
    \[
      \sum_{k=1}^{n} k = \frac{n(n+1)}{2}
    \]
  }{%
    \textbf{Практическая часть}\\[6pt]
    \begin{itemize}
      \item Эксперимент~A
      \item Эксперимент~B
      \item Эксперимент~C
    \end{itemize}
  }
\end{frame}

% [nonavbar] — слайд не регистрируется в навигационных точках
\begin{frame}[nonavbar]{Резервный слайд (без точки в шапке)}
  \framesubtitle{Опция \texttt{[nonavbar]}}
  Этот фрейм не получает точку в навигационной полосе.
  Подходит для резервных, вспомогательных или переходных слайдов,
  которые не должны загромождать навигацию.
\end{frame}

% ── Section divider ──────────────────────────────────────────
\sectionslide{Заключение}

% ============================================================
\section{Заключение}
% ============================================================

\begin{frame}{Выводы}
  \begin{emphblock}{Главный результат}
    Краткая формулировка основного итога работы,
    видимая с последнего ряда аудитории.
  \end{emphblock}

  \begin{enumerate}
    \item Первый подтверждённый вывод
    \item Второй вывод с количественной оценкой
    \item Направления дальнейших исследований
  \end{enumerate}
\end{frame}

% ── FINAL SLIDE ──────────────────────────────────────────────
\begin{frame}[plain,noframenumbering]
  \finalslide{\ContactURL}{\FeedbackURL}
\end{frame}

\end{document}
